\section{Output files\label{sec:Output-files}}

A series of specific directives causes the writing of output files
(see Tab \ref{tab:inputlist}). In JETSPIN three output files can
be written: 1) the file \texttt{statdat.dat} containing time-dependent
statistical data of simulated process; 2) the file \texttt{traj.xyz}
reporting the nanofiber trajectory in XYZ format file; 3) the files
\texttt{frame\%06d.xyz} reporting the nanofiber geometries in splitted
XYZ format files (a single file for each frame); 4) the files \texttt{frame\%06d.pdb}
reporting the nanofiber geometries in splitted PDB (Protein Data Bank)
format files (a single file for each frame). Note that in the last
case a topological file \texttt{frame\%06d.psf} in PSF (Protein Structure
File) format will be printed alongside each pdb file. These two files
(PDB and PSF) can be opened by VMD-Visual Molecular Dyanamics \cite{humphrey1996vmd}.

Various statistical data can be written on the file \texttt{statdat.dat}
, and the user can select them in input using the directive \texttt{printstat
list} followed by appropriate symbolic strings, whose the list with
corresponding meanings is reported in Tab \ref{tab:outputlist}. The
selected observables will be printed at time interval on the same
line following the order specified in input file. In the same way,
a list of statistical data can be printed on computer terminal using
the directive \texttt{print list} followed by the symbolic strings
of Tab \ref{tab:outputlist}.

At normal termination, the terminal also reports the wall-clock duration of
the temporal integration loop and its throughput in integration steps per
second. The timer starts immediately before the loop and stops immediately
after it; input parsing, initial allocation, restart loading, and output-file
setup are therefore excluded. Integration, bead insertion and removal,
statistics, and scheduled output executed inside the loop are included. In
MPI builds, ranks are synchronized at both boundaries and rank zero reports
the elapsed duration of the slowest participating rank. This loop-only wall
time is the recommended measurement for CPU, MPI, and accelerator performance
comparisons.

For performance diagnosis, setting the environment variable
\texttt{JETSPIN\_PROFILE=1} prints an optional wall-clock breakdown for the
integrator, Coulomb calculation, bead topology operations, statistics,
scheduled output, and restart output. Profiling is disabled by default. The
Coulomb value is a nested part of the integrator total and must not be added
to it when interpreting percentages. The profiler prints cumulative wall
time, percentage of the complete temporal-loop time, and call count on the
terminal; it does not write a profiling file. These are explicitly
instrumented regions rather than an automatic measurement of every Fortran
subroutine. Region timings are inclusive, so the integrator value includes
the Coulomb calls.

In MPI execution, the detailed region breakdown is local to rank zero and is
not reduced to the maximum over all ranks. The synchronized complete-loop
wall time should be used for MPI scaling. Since clock reads add a small
overhead, performance comparisons should enable profiling in both runs or
disable it in both runs.

The file \texttt{traj.xyz} is written as a continuous series of XYZ
format frames taken at time interval, so that can be read by suitable
visualization programs (e.g. VMD-Visual Molecular Dynamics \cite{humphrey1996vmd},
UCSF Chimera \cite{pettersen2004ucsf}, etc.) to generate animations.
The number of elements contained in the file is kept constant equal
to the value specified by the directive \texttt{print xyz maxnum},
since few programs (e.g. VMD) do not manage a variable number of elements
along the simulations. If the actual number of beads is lower than
the given constant \texttt{maxnum}, the extra elements are printed
in the origin point. On the other hand, the number of beads printed
in the files \texttt{frame\%06d.xyz} and \texttt{frame\%06d.pdb} is
exactly the actual number of beads used to discretized the jet at
the given frame. Note that the unit of mesurement used in the files,
\texttt{traj.xyz}, \texttt{frame\%06d.xyz} and \texttt{frame\%06d.pdb}
is the centimeter multiplied by $f$, which was given in input file
by the directve `` \textit{print xyz rescale} $f$ '' and `` \textit{print
pdb rescalepdb} $f$ '' (see Table \ref{tab:inputlist}), respectively.

\begin{longtable}{cc}
\caption{In JETSPIN a series of instantaneous and statistical data are available
to be printed by selecting the appropriate key. Here, the list of
symbolic string keys is reported with their corresponding meanings.
Note by \textit{scaled} we mean that the observable was rescaled by
the characteristic values provided in Tab\ref{tab:tabella-adimensionale}.
The underlined keys correspond to data which are averaged on the time
interval given by the directive \texttt{print time} in input file.
Note by \textit{farest bead} we mean the farest bead from the nozzle
(with greatest $x$ value). Note all the quantities are expressed
in centimeter\textendash gram\textendash second unit system, excepted
the dimensionless scaled observables.}
\tabularnewline
\endlastfoot
\hline
\textbf{keys:}  & \textbf{meaning:}\tabularnewline
\hline
\hline
t  & unscaled time\tabularnewline
ts  & scaled time \tabularnewline
x, y, z  & unscaled coordinates of the farest bead \tabularnewline
xs, ys, zs  & scaled coordinates of the farest bead\tabularnewline
st  & unscaled stress of the farest bead\tabularnewline
sts  & scaled stress of the farest bead\tabularnewline
vx, vy, vz  & unscaled velocities of the farest bead \tabularnewline
vxs, vys, vzs  & scaled velocities of the farest bead \tabularnewline
yz  & unscaled normal distance from the x axis of the farest bead\tabularnewline
yzs  & scaled normal distance from the x axis of the farest bead\tabularnewline
mass  & last inserted mass at the nozzle \tabularnewline
q  & last inserted charge at the nozzle\tabularnewline
cpu  & time for every print interval\tabularnewline
cpur  & remaining time to the end \tabularnewline
cpue  & elapsed time \tabularnewline
n  & number of beads used to discretise the jet\tabularnewline
f  & index of the first bead \tabularnewline
l  & index of the last bead \tabularnewline
\uline{curn}  & current at the nozzle \tabularnewline
\uline{curc}  & current at the collector \tabularnewline
\uline{vn}  & velocity modulus of jet at the nozzle \tabularnewline
\uline{vc}  & velocity modulus of jet at the collector\tabularnewline
\uline{svc}  & strain velocity at the collector \tabularnewline
\uline{mfn}  & mass flux at the nozzle \tabularnewline
\uline{mfc}  & mass flux at the collector \tabularnewline
\uline{rc}  & radius of jet at the collector \tabularnewline
\uline{rrr}  & radius reduction ratio of jet \tabularnewline
\uline{lp}  & length path of jet \tabularnewline
\uline{rlp}  & length path of jet divided by the collector distance\tabularnewline
mxst  & maximum stress along the jet\tabularnewline
mxsx  & x position of the jet maximum stress\tabularnewline
\hline
\label{tab:outputlist}  & \tabularnewline
\end{longtable}
